\documentclass[10pt,a4paper]{article}
\usepackage[T1]{fontenc}
\usepackage{newtxtext,newtxmath}
\usepackage[a4paper,top=2.2cm,bottom=2.2cm,left=2.3cm,right=2.3cm]{geometry}
\usepackage{xcolor}
\usepackage[skins]{tcolorbox}
\usepackage{booktabs}
\usepackage{tabularx}
\usepackage{enumitem}
\usepackage{url}
\usepackage{titlesec}

\setlength{\emergencystretch}{1.5em}
\setlength{\parindent}{0pt}
\setlength{\parskip}{4pt}
\setlist{nosep,leftmargin=1.3em,topsep=2pt}
\titleformat{\section}{\large\bfseries}{}{0pt}{}
\titlespacing*{\section}{0pt}{12pt}{4pt}
\titleformat{\subsection}{\normalsize\bfseries}{}{0pt}{}
\titlespacing*{\subsection}{0pt}{8pt}{2pt}

\definecolor{cmtbg}{gray}{0.94}
\definecolor{cmtrule}{gray}{0.55}
\newtcolorbox{cmt}{enhanced,sharp corners,boxrule=0pt,colback=cmtbg,
  left=6pt,right=6pt,top=3pt,bottom=3pt,before skip=4pt,after skip=4pt,
  borderline west={2pt}{0pt}{cmtrule},fontupper=\itshape}
\newcommand{\resp}{\textbf{Response.}\ }
\newcommand{\changes}[1]{\par\textbf{Changes in manuscript:} #1}
\newcommand{\pH}{p_{\mathrm{H}}}
\newcommand{\Fall}{F_{\text{all}}}
\newcommand{\pp}{\,pp}

\begin{document}

\begin{center}
{\LARGE\bfseries Response to Reviewers}\\[6pt]
{\large ICIC--APTIKOM 2026, Paper ID 57}\\[4pt]
\textit{Benchmarking Multi-View Tree-Level Oil Palm Bunch Counting Under a Fixed Detector}
\end{center}

\vspace{2pt}
\begin{tabularx}{\textwidth}{@{}lX@{}}
Authors & Muhammad Zainal Muttaqin, Fatma Indriani, Setyo Wahyu Saputro, Alia Rahmi,
          Triando Hamonangan Saragih, Hartoni, Dwi Kartini, Naufal Said\\
Corresponding author & Muhammad Zainal Muttaqin (mz.muttaqin1@gmail.com)\\
Revised manuscript & Camera-ready version, 6 pages, certified by IEEE PDF eXpress (file ID 2026332983)\\
\end{tabularx}

\section{General Response}
We thank the four reviewers for their careful reading and constructive comments. We
addressed every comment in the revised manuscript. The revision keeps the original
experimental design and the six-page limit, and adds the following analyses:
\begin{itemize}
\item Tree-level bootstrap 95\% confidence intervals (CIs), exact paired permutation
  tests and Holm correction for all counter and feature comparisons
  (Sections II-E and III-A to III-C).
\item A second detector family (YOLO11m) trained on the official split, four YOLO26
  checkpoints and a confidence-threshold sweep (Sections II-B and III-E, Table VI).
\item Appearance-level matching that separates missed annotations, wrong-class
  detections and false positives, with diagnostic counting rules
  (Table I, Section III-D, Fig.~2).
\item Per-class MAE, RMSE, signed bias and total-count error
  (Section II-D, Tables III and V).
\item Validation-based model selection, matched repeated cross-validation (CV) and
  cross-estate counter transfer (Sections II-E, III-C and III-E).
\item A feature-definition table (Table II), an analysis of the Random Forest (RF)
  results (Section III-F) and a limitations subsection (Section III-G).
\end{itemize}
During the re-analysis we also found and corrected three errors in the submitted
version; they are listed in Response 3.10.

Section, table and figure numbers below refer to the revised manuscript. The submitted
Tables II and III are merged into the revised Table III, the submitted Table V is now
Table IV, and the submitted Table VI is now Table V. Code, cached predictions and the
outputs of all added analyses are available at
\url{https://github.com/ULM-SawitMVC/Baseline-SawitMVC} (folders
\texttt{experiments/revision} and \texttt{results/revision}).
The abbreviation pp denotes percentage points, and $\pH$ denotes a Holm-adjusted
$p$-value.

%======================================================================
\section{Reviewer 1}

\begin{cmt}
CONCLUSION: Include further implications for the field. Discuss benefits or
shortcomings of your work and suggest future areas for research.
\end{cmt}
\resp We thank the reviewer for the positive assessment. We rewrote Section IV to cover
the three requested points.
\begin{itemize}
\item \emph{Implications.} Evaluating the same counter under ground-truth (GT) and
  detector inputs lets practitioners separate the effect of input quality from the
  effect of aggregation. For the Black Bunch Census, the results support improving B3
  and B4 detection and B2-to-B3 classification together with aggregation. Class-level
  and total-count metrics should both be reported, because a plausible total count can
  hide class errors that change the harvest schedule.
\item \emph{Benefits.} Cached predictions, fixed splits and released code allow new
  counters to be compared without repeating detector inference.
\item \emph{Shortcomings.} The evidence covers one dataset from two estates of unequal
  size, and acquisition changes are untested. The results do not give an
  association-error budget or a limit on what future counters can achieve.
  Section III-G lists the remaining limitations.
\item \emph{Future research.} Comparison of explicit multi-view association with learned
  aggregation under matched inputs, broader detector evaluation, repeated detector
  training, and tests on an independent plantation and at measured camera-pitch angles.
\end{itemize}
\changes{Section IV (Conclusion); Section III-G (Limitations).}

%======================================================================
\section{Reviewer 2}

\begin{cmt}
However, the current manuscript does not yet provide sufficient evidence to support some
of its broader conclusions. The use of a single detector, limited test set, absence of
statistical significance analysis, and limited error decomposition are the main
weaknesses. I would recommend acceptance after minor revision, provided that the authors
strengthen the statistical analysis, clarify the scope of their conclusions, and provide
a more rigorous decomposition of detector and counting errors.
\end{cmt}
\resp We thank the reviewer for recognising the value of the controlled comparison. We
addressed the four weaknesses as follows. The responses to Reviewer~3 give further
details.
\begin{enumerate}[label=(\alph*)]
\item \emph{Single detector.} We trained a second detector family, YOLO11m, on the
  official 716/96/141 split with the same training settings. With the same Ridge$+F_0$
  counter on the 141 test trees, Class $\pm$1 accuracy is 73.76\% for YOLO11m, 76.06\%
  for YOLO26m and 97.70\% for GT inputs. GT inputs remain 23.94\pp{} better than YOLO11m
  (95\% CI 20.57--27.30; $p<0.001$), while the two detectors do not differ significantly
  ($-2.30$\pp; CI $-4.96$ to $+0.35$; $p=0.118$). Four YOLO26 checkpoints span 3.13\pp,
  compared with a 23.44\pp{} gap to GT inputs (Response 3.2).
\item \emph{Limited test set.} Matched repeated CV uses identical folds under both input
  conditions. On 237 trees outside detector gradient training, Class $\pm$1 is
  $97.72\pm1.02$\% (GT) versus $72.70\pm2.35$\% (fixed detector); on the 141 original
  test trees it is $97.16\pm1.23$\% versus $74.08\pm2.96$\%. A validation-only model
  selection gives a gap of 23.58\pp{} (Responses 3.1 and 3.5).
\item \emph{Significance analysis.} The class-level gap between the best configurations
  is 20.57\pp{} (CI 17.20--23.94; $p<0.001$). With the same Ridge$+F_0$ counter in both
  conditions, it is 21.63\pp{} (CI 18.44--25.00; $p<0.001$), so the gap is not caused
  by comparing different counter families. No pair of counters differs significantly
  under the fixed detector after Holm correction (Response 3.1).
\item \emph{Error decomposition.} Appearance-level matching reports missed, wrong-class
  and false-positive counts for each class, and 22.7\% of test bunches are not detected
  in any view. Diagnostic counting rules show the effect of each error type
  (Responses 3.3 and 3.9).
\end{enumerate}
We also narrowed the conclusions to the tested configurations on SawitMVC
(Response 3.8).
\changes{Abstract; Sections II-B, II-E, III-A to III-E and IV; Tables I and VI; Fig.~2.}

%======================================================================
\section{Reviewer 3}

\subsection{Comment 3.1}
\begin{cmt}
Add statistical uncertainty/significance analysis for the reported model differences.
\end{cmt}
\resp Section II-E defines the protocol. Percentile bootstrap CIs resample trees 10,000
times (seed 42), keep the four class counts of each tree together, and reuse the same
resamples for paired comparisons. Exact paired permutation tests exchange the outcomes of
two models within each tree. Holm correction covers the ten counter pairs within each
input condition and, as a separate family, the seven feature sets compared with $F_0$.
Wilcoxon tests on per-tree macro absolute error and exact McNemar tests on Tree $\pm$1
cover the other endpoints. The main results are:
\begin{itemize}
\item GT inputs: the four counters other than RF do not differ significantly ($\pH=1$).
  ElasticNet exceeds RF by 2.13\pp{} (CI 0.89--3.55; $p=0.0042$, $\pH=0.0376$), and SVM
  also exceeds RF ($\pH=0.0342$).
\item Fixed detector: none of the ten counter pairs differs significantly. The largest
  difference, ElasticNet versus RF, is 3.19\pp{} ($p=0.0628$, $\pH=0.628$).
\item Input conditions: 20.57\pp{} (CI 17.20--23.94) between the best configurations and
  21.63\pp{} (CI 18.44--25.00) for the same counter; both $p<0.001$.
\end{itemize}
The manuscript states that a nonsignificant difference does not establish equivalence.
Because the submitted version highlighted the configuration with the best test score,
we added a validation-only selection over 40 model and feature candidates per condition.
The selected configurations, Ridge$+F_0$ (GT) and SVM$+(F_0+\text{conf})$ (fixed),
reach 97.70\% and 74.11\% on the test trees, a gap of 23.58\pp{} (CI 20.21--27.13;
$p<0.001$).
\changes{Sections II-E, III-A, III-B and III-C.}

\subsection{Comment 3.2}
\begin{cmt}
Clarify the scope of the detector-related conclusion, because only one detector
checkpoint is evaluated.
\end{cmt}
\resp We agree that one checkpoint cannot support a general statement. We added three
analyses (Table VI) and restricted the conclusion accordingly.
\begin{itemize}
\item \emph{Detector family.} YOLO11m was trained on the official split with 60 epochs,
  batch 32, image size 640, patience 60 and seed 42; its checkpoint was selected on the
  validation trees. The results are given in Response 2(a). YOLO11m misses 976 and
  misclassifies 399 of the 2,612 GT appearances.
\item \emph{Checkpoints.} YOLO26n, YOLO26s and YOLO26m from an earlier split, and YOLO26m
  from the official split, reach 70.70--73.83\% Class $\pm$1 on 64 trees outside the
  gradient training of all checkpoints, compared with 97.27\% for GT inputs.
\item \emph{Operating point.} Over eleven confidence thresholds from 0.05 to 0.70, with
  the counter refitted at each threshold, no threshold exceeds 77.48\% Class $\pm$1.
  Threshold 0.10 raises Tree $\pm$1 from 32.62\% to 38.30\%.
\end{itemize}
Section III-G states the remaining limits: one training run per detector, different
Ultralytics versions and defaults (8.4.49 and 8.4.140), and confounding of model
capacity with training split in the checkpoint analysis. The manuscript therefore
describes the finding as applying to the tested detectors on SawitMVC, not as a general
ranking of architectures.
\changes{Sections II-B, III-E and III-G; Table VI.}

\subsection{Comment 3.3}
\begin{cmt}
Provide more detailed error decomposition for missed detections, false positives, and
maturity-class confusion.
\end{cmt}
\resp Section III-D matches every detection on the test trees to the annotations. The
matching is greedy in order of confidence, class-agnostic, and accepts a pair at
IoU $\geq 0.5$. Table I reports, for each class, annotations with no matched detection
(Missed), matched boxes with an incorrect class (Wrong) and detections with no matched
annotation (FP).
\begin{itemize}
\item Of 2,612 annotated appearances, 964 (36.9\%) are missed, 396 (15.2\%) have an
  incorrect class and 1,252 (47.9\%) are correct.
\item Of 2,354 detections, 706 (30.0\%) are false positives. Their mean confidence is
  0.337, compared with 0.487 for matched detections.
\item B2 errors are mainly class confusion (209 appearances, of which 158 are predicted
  as B3, versus 161 missed). B4 errors are mainly missed detections (262 of 455).
\item Of 1,397 physical test bunches, 317 (22.7\%) are not detected in any view, from
  8.5\% for B1 to 42.3\% for B4.
\end{itemize}
The manuscript notes that an unmatched annotation can also result from an IoU below the
threshold or from matching competition.
\changes{Table I; Section III-D.}

\subsection{Comment 3.4}
\begin{cmt}
Add additional counting metrics, especially per-class MAE/RMSE and total bunch-count
error.
\end{cmt}
\resp Section II-D defines macro MAE, macro RMSE, signed class bias, total-count MAE and
total $\pm$1 accuracy. Table III reports macro MAE for all counters under both input
conditions, and Table V gives per-class accuracy, MAE, RMSE and bias for the best
configuration in each condition. From GT to fixed-detector inputs, macro RMSE rises
from 0.522 to 1.418, total-count MAE rises from 0.979 to 1.787, and total $\pm$1
accuracy falls from 78.01\% to 48.94\% (Section III-C). The manuscript distinguishes
total $\pm$1, in which class errors can cancel, from Tree $\pm$1, which requires all
four class counts to be within one bunch.
\changes{Sections II-D and III-C; Tables III and V.}

\subsection{Comment 3.5}
\begin{cmt}
Strengthen validation using repeated splits, cross-validation, or an independent
plantation if available.
\end{cmt}
\resp No independent plantation was available for this revision. We therefore added
three resampling analyses and state the missing independent test as a limitation.
\begin{itemize}
\item \emph{Matched repeated CV.} Five folds with five repeats, stratified by estate and
  dominant class, use identical folds for GT and fixed-detector inputs with the
  detector frozen. Class $\pm$1 is $97.72\pm1.02$\% versus $72.70\pm2.35$\% on 237 trees
  outside detector gradient training (gap 25.03\pp), and $97.16\pm1.23$\% versus
  $74.08\pm2.96$\% on the 141 original test trees (gap 23.08\pp).
\item \emph{Validation-based selection.} Model and features are selected on the
  validation trees only; the test gap is 23.58\pp{} (Response 3.1).
\item \emph{Cross-estate transfer.} Ridge$+\Fall$ trained on 641 DAMIMAS trees reaches
  71.88\% on 24 LONSUM evaluation trees, compared with 72.92\% when trained on 75
  LONSUM trees. The reverse transfer reaches 68.78\% on 213 DAMIMAS trees. Because the
  detector was trained on both estates, this analysis evaluates the counter only.
\end{itemize}
\changes{Sections II-E, III-C, III-E and III-G.}

\subsection{Comment 3.6}
\begin{cmt}
Expand the feature-definition table to improve reproducibility.
\end{cmt}
\resp Table II now defines all 67 dimensions: $F_0$ (13), confidence (20), spatial (8),
distribution (20) and composition (6). It gives the formulas, the normalisation of box
position and area by image size, the confidence thresholds ($p\geq0.5$ and $p\geq0.6$),
the population standard deviation over all views including empty views,
$\varepsilon=10^{-6}$, and the default values for a class with no detection. Section II-C
defines the notation, and the released feature extractor produces the same 67
dimensions. We also report that the 1.42\pp{} gain of $\Fall$ over $F_0$ is not
statistically resolved (CI $-0.53$ to $+3.37$; $p=0.201$, $\pH=1$) and that $\Fall$ does
not lead under repeated CV (Table IV).
\changes{Section II-C; Tables II and IV.}

\subsection{Comment 3.7}
\begin{cmt}
Discuss the poor Random Forest performance more thoroughly.
\end{cmt}
\resp Section III-F explains the result. Each regression tree averages the training
targets within a leaf, so RF cannot predict beyond the training count range and pulls
estimates towards common counts when high counts are rare. The diagnostics agree with
this explanation:
\begin{itemize}
\item Slopes of predicted on true counts are 0.754--0.888 for RF and 0.865--0.951 for
  the linear counters, and the relative prediction standard deviation of RF is smaller
  in every class.
\item RF never predicts more than three B1 bunches, although the maximum is five in
  training and six in the test set; its B1 MAE is 0.149, compared with 0.050 for Ridge.
\item In the highest-load bin, RF trails ElasticNet by 3.07\pp{} (94.74\% versus 97.81\%).
\end{itemize}
Under GT inputs, RF is significantly worse than ElasticNet and SVM after Holm correction;
under the fixed detector the difference is not resolved ($p=0.0628$, $\pH=0.628$). The
manuscript presents these findings as diagnostics and does not exclude improvement by
tuning.
\changes{Section III-F.}

\subsection{Comment 3.8}
\begin{cmt}
Moderate broad generalizations about detector quality versus counter quality.
\end{cmt}
\resp We revised the following statements:
\begin{itemize}
\item The submitted abstract stated that ``the dominant finding is that detector output
  quality, more than the choice of counter, limits deployed accuracy.'' The revised
  abstract states that the results ``support prioritising detector output quality on
  this dataset'' and ``do not establish a ceiling on counter improvement.''
\item The submitted statement that the remaining counting problem is ``simple enough for
  compact regression models'' was removed.
\item GT inputs are described as an empirical reference, not a mathematical upper bound
  on future counting methods (Section II-B).
\item Nonsignificant differences are not interpreted as equivalence (Sections II-E, III-A
  and III-E).
\item The conclusion is limited to the tested configurations and lists the evidence
  needed before the claims can be extended to operational harvest forecasts
  (Section IV).
\end{itemize}
\changes{Abstract; Sections II-B, II-E, III-A to III-G and IV.}

\subsection{Comment 3.9}
\begin{cmt}
Clearly distinguish detector performance limitations from multi-view
association/aggregation limitations.
\end{cmt}
\resp Section III-D and Fig.~2(a) compare six counting rules on the 141 test trees:

\begin{center}
\small
\begin{tabular}{@{}lr@{}}
\toprule
Counting rule & Class $\pm$1 (\%) \\
\midrule
GT boxes and classes, ElasticNet$+F_0$ & 98.05 \\
Matched physical bunches counted once, GT class & 86.17 \\
Same identities, confidence-weighted detector class vote & 71.81 \\
Class vote, each unmatched detection counted as a separate bunch & 62.23 \\
Deployed counter, Ridge$+\Fall$ & 77.48 \\
Naive sum of detections & 50.71 \\
\bottomrule
\end{tabular}
\end{center}

The three identity-based rules show the separate effects of bunches that are never
detected, class confusion and unmatched detections. The manuscript states that these
rules are separate estimators and that their differences do not add up to the deployed
error. The deployed counter (77.48\%) exceeds two of the rules because learned
aggregation compensates statistically for missed counts and class bias. A direct
measurement of association error requires an explicit association module, which we list
as future work.
\changes{Section III-D; Fig.~2(a); Section IV.}

\subsection{Comment 3.10}
\begin{cmt}
Carefully proofread the manuscript for formatting, equations, tables, and citation
consistency.
\end{cmt}
\resp We proofread the full manuscript. All tables were re-typeset (Response 4.1),
equations (1) and (2) define Class $\pm$1 and Tree $\pm$1, abbreviations are defined at
first use, and references are numbered in order of first citation. The final PDF passed
IEEE PDF eXpress. During the re-analysis we found and corrected the following errors in
the submitted version:
\begin{itemize}
\item \emph{Global divisor (GT inputs).} Class $\pm$1 / Tree $\pm$1 / MAE were corrected
  from 95.39\% / 85.11\% / 0.376 to 95.57\% / 86.52\% / 0.356. The divisor is fitted on
  the 716 training trees only, separately for each input condition ($k=1.891$ for GT
  and $1.793$ for the fixed detector). A fixed-detector row (70.21\% / 22.70\% / 1.188)
  was added to Table III.
\item \emph{Tree $\pm$1 gap.} The value was corrected from 59.58 to 59.57\pp.
\item \emph{Visibility-adaptive divisor.} This row was removed because we could not
  confirm that the trees used to develop the heuristic are disjoint from the current
  test split.
\item \emph{Software citation.} The unsupported version number ``26.0.0'' was removed
  from the Ultralytics reference; the training log records Ultralytics 8.4.49.
\end{itemize}
\changes{All sections; Tables I to VI; references.}

%======================================================================
\section{Reviewer 4}

We thank the reviewer for the positive assessment of the problem statement, dataset and
methodology.

\subsection{Comment 4.1}
\begin{cmt}
Fix several noticeable typesetting and column-alignment errors in Tables III and V.
\end{cmt}
\resp The submitted Table III (fixed-detector counting) is now merged with the submitted
Table II into the revised Table III, which places GT and fixed-detector results side by
side. The submitted Table V (feature ablation) is now Table IV. All six tables were
re-typeset with 8-point text, horizontal rules only, the same number of decimal places
in each column, and widths within the column. Captions were shortened, and units and
definitions were moved to table notes.
\changes{Tables I to VI.}

\subsection{Comment 4.2}
\begin{cmt}
Briefly evaluate an additional detector architecture to strengthen the generalizability
of the findings.
\end{cmt}
\resp We trained YOLO11m on the official 716/96/141 split with the same epochs, batch
size, image size, patience and seed as the YOLO26m detector, and selected its checkpoint
on the validation trees. With the same Ridge$+F_0$ counter on the 141 test trees:

\begin{center}
\small
\begin{tabular}{@{}lccc@{}}
\toprule
Input & Detections & Class $\pm$1 (\%) & Tree $\pm$1 (\%) \\
\midrule
YOLO11m & 2,099 & 73.76 & 26.24 \\
YOLO26m & 2,354 & 76.06 & 28.37 \\
GT detections & 2,612 & 97.70 & 90.78 \\
\bottomrule
\end{tabular}
\end{center}

The two detectors do not differ significantly in Class $\pm$1 ($-2.30$\pp; CI $-4.96$ to
$+0.35$; $p=0.118$), and GT inputs remain 23.94\pp{} better than YOLO11m (CI
20.57--27.30; $p<0.001$). The manuscript notes that one training run per detector and
different software versions limit wider generalisation.
\changes{Sections II-B, III-E and III-G; Table VI.}

\subsection{Comment 4.3}
\begin{cmt}
Clarify how robust the vertical spatial features are to real-world camera pitch
variations.
\end{cmt}
\resp The dataset has no calibrated pitch measurements and no repeated captures at known
angles, so physical pitch robustness could not be tested. We added a sensitivity analysis
of the vertical feature. The mean vertical position of the test features was transformed
as $\overline{cy}\mapsto\operatorname{clip}(a\,\overline{cy}+b,0,1)$ over 35 settings of
$a\in[0.90,1.10]$ and $b\in[-0.10,0.10]$, with training data unchanged and missing-class
values kept at 0.5. Ridge$+\Fall$ loses at most 1.42\pp{} Class $\pm$1, and Gaussian noise
with $\sigma=0.05$ costs 0.53\pp. This analysis tests feature shifts only. A real pitch
change also affects perspective, box area, visibility and detector output. The manuscript
therefore states that robustness to camera rotation remains unverified
(Sections III-E and III-G) and proposes repeated captures at measured pitch angles as
future work (Section IV).
\changes{Sections III-E, III-G and IV.}

\end{document}
